% vim: spelllang=en spell


\documentclass[12pt,notitlepage]{report}
%\pagestyle{headings}
\pagestyle{plain}

%\frenchspacing

\usepackage[utf8]{inputenc}
%\usepackage{czech}
\usepackage{a4wide}
%\usepackage{index} % nutno použít v případě tvorby rejstříku balíčkem makeindex
%\usepackage{fancybox} % umožňuje pokročilé rámečkování :-)
\usepackage{graphicx} % nezbytné pro standardní vkládání obrázků do dokumentu

\usepackage[left=4cm]{geometry}
\usepackage{xspace}

%\newindex{default}{idx}{ind}{Rejstřík} % zavádí rejstřík v případě použití balíku index

\title{Trojita: A Qt IMAP Client}
\author{Jan Kundrát}

\newcommand{\trojita}{Trojita\xspace}

\hyphenation{WebKit Imap Trojita}

%\date{}

\begin{document}


\begin{titlepage}
\begin{center}
%\ \\

\vspace{15mm}

\large
Charles University in Prague\\
Faculty of Mathematics and Physics\\

\vspace{5mm}

{\Large\bf BACHELOR THESIS}

\vspace{10mm}

\includegraphics[scale=0.3]{logo.eps}

\vspace{15mm}

%\normalsize
{\Large Jan Kundrát}\\
\vspace{5mm}
{\Large\bf IMAP E-mail Client}\\
\vspace{5mm}
Department of Software Engineering\\
%\end{center}
\vspace{15mm}

\large
\noindent Supervisor: Mgr. Vlastimil Babka
\vspace{1mm} 

\noindent Study Program: Computer Science, Programming

\vspace{20mm}

%\begin{center}
2009
\end{center}

\end{titlepage}

\normalsize
\setcounter{page}{2}
\ \vspace{10mm} 

\noindent I'd like to thank my supervisor, Mgr. Vlastimil Babka, for his numerous
advices during the writing of this thesis, and Ms. Anna Adamcová for her great
patience and support.

\vspace{\fill}
\noindent I hereby declare that I wrote this thesis myself using only the
referenced sources. I also agree with lending and publishing of this thesis.

\bigskip
\noindent Prague, May 29, 2009 \hspace{\fill}Jan Kundrát\\

\tableofcontents

\newpage

\noindent
Název práce: IMAP E-mail Client\\
Autor: Jan Kundrát\\
Katedra (ústav): Katedra softwarového inženýrství\\
Vedoucí bakalářské práce: Mgr. Vlastimil Babka\\
e-mail vedoucího: Vlastimil.Babka@mff.cuni.cz\\

\noindent Abstrakt:  V předložené práci studujeme ... Uvede se abstrakt v
rozsahu 80 až 200 slov. Lorem ipsum dolor sit amet, consectetuer adipiscing
elit. Ut sit amet sem. Mauris nec turpis ac sem mollis pretium. Suspendisse
neque massa, suscipit id, dictum in, porta at, quam. Nunc suscipit, pede vel
elementum pretium, nisl urna sodales velit, sit amet auctor elit quam id tellus.
Nullam sollicitudin. Donec hendrerit. Aliquam ac nibh. Vivamus mi. Sed felis.
Proin pretium elit in neque. Pellentesque at turpis. Maecenas convallis.
Vestibulum id lectus. Fusce dictum augue ut nibh. Etiam non urna nec mi mattis
volutpat. Curabitur in tortor at magna nonummy gravida. Mauris turpis quam,
volutpat quis, porttitor ut, condimentum sit amet, felis.\\

\noindent Klíčová slova: IMAP, e-mail, SMTP

\vspace{10mm}

\noindent
Title: IMAP E-mail Client\\
Author: Jan Kundrát\\
Department: Department of Software Engineering\\
Supervisor: Mgr. Vlastimil Babka\\
Supervisor's e-mail address: Vlastimil.Babka@mff.cuni.cz\\

\noindent Abstract: In the present work we study ... Uvede se anglický abstrakt
v rozsahu 80 až 200 slov. Lorem ipsum dolor sit amet, consectetuer adipiscing
elit. Ut sit amet sem. Mauris nec turpis ac sem mollis pretium. Suspendisse
neque massa, suscipit id, dictum in, porta at, quam. Nunc suscipit, pede vel
elementum pretium, nisl urna sodales velit, sit amet auctor elit quam id tellus.
Nullam sollicitudin. Donec hendrerit. Aliquam ac nibh. Vivamus mi. Sed felis.
Proin pretium elit in neque. Pellentesque at turpis. Maecenas convallis.
Vestibulum id lectus. Fusce dictum augue ut nibh. Etiam non urna nec mi mattis
volutpat. Curabitur in tortor at magna nonummy gravida. Mauris turpis quam,
volutpat quis, porttitor ut, condimentum sit amet, felis. \\

\noindent Keywords: IMAP, e-mail, SMTP

\newpage

\chapter{Introduction}

Hic sunt Leones.

\subsection{Motivation}

While there are certainly numerous stable and widely-used Mail User Agents on
the market, none of them fulfils my expectations of reasonable performance and
efficiency.  Indeed, most of these applications started as MUAs speaking the
POP3 protocol \footnote{The POP3 protocol\ldots FIXME}, adding support for IMAP
later on.  Typically, this support was added long since the initial design phase
of the product was completed, resulting in suboptimal design
choices.\cite{crispin-ten-commandments} Some clients, like the \ldots FIXME

\subsection{Structure of the thesis}

In the following chapter, we provide a gentle introduction to the world of
IMAP and related standards like MIME.  Following that, we move to explain some
design choices made during the development proces of \trojita, as well as
elaborate on the internal application architecture and structure.  User
documentation is the next chapter, and we finish the whole thesis by a wrap-up
providing an overview of what we implemented and how is the result usable.

\chapter{IMAP and Related Technologies}

Electronic mail, or a SMTP-mail, is a public service suitable for automated
message exchange among connected entities.  Interoperability is guaranteed by
several {\em Internet standards}, usually codified in the form of RFC documents.
Subject to these standards are various protocols specifying the rules of what
entities can communicate to each other as well as definition of the format of
all transmitted messages. In this section, we provide a gentle introduction to
the numerous standards which deal with this highly complex topic.

\section{Basic Concepts}

An {\em email message}, or a RFC-822\footnote{As defined by the RFC 822 standard
and subsequently modified by RFC 2822 \cite{rfc-2822} and others} message,
consists of three parts: {\em Envelope}, {\em Header} and {\em Body}. The {\em
envelope} is the only relevant part for mail exchange among Internet
hosts; it includes basic information like addresses of the sender and recipient.
A {\em header} provides more detailed metadata about the message, from
human-readable sender addresses and message route traces to user-defined
fields.  Some pieces of this information are used when the original delivery
fails for some reason, like the {\tt Return-Path} header.  The {\em header} also
serves as a basis for an {\em IMAP envelope data} which is explained later
in chapter \ref{imap-envelope}.  Finally, a {\em message body} is what an
average user typically refers to as ``the e-mail''.  It might contain just a
plain US-ASCII text, an HTML message with embedded images or it could be a
recursively defined entity with rich tree structure.

An {\em IMAP server} is a host in the Internet which provides access to local
mail store via the standard protocol, IMAP4rev1 in this case. The IMAP server
might be located in an employer's data center or on the user's laptop, for
example.

An {\em MTA}\footnote{Mail Transfer Agent} or an {\em SMTP Server} is an
Internet host whose purpose is delivery of the RFC 2822 mail.  MTAs in the
Internet communicate with each other, determining routing information from
MX-records\footnote{Mail eXchange, a special record in the Domain Name Service
which specifies what servers are responsible for accepting all incoming e-mail
for the particular domain} in the DNS.  As a common practice nowadays, these
servers also accept outgoing e-mail from their own users, provided the
connection is properly authenticated, either by simple fact that it originates
from a trusted network or that the user has provided her credentials, typically as
a user name and password combination or in the form of an X.509 certificate.

A {\em Mailbox} or a {\em Mail Folder} is a place on the IMAP server where the
messages are logically stored.  A mailbox might contain number of child
mailboxes if the server implementation allows it.  Certain IMAP software has to
be told in advance whether the mailbox just being created is supposed to contain
only other mailboxes or only regular e-mail messages, while others allow free
mixing of both.

\section{IMAP-specific Attributes}

A message in the IMAP sense is {\em immutable}, that is, its contents (like
header, structure of the body or the actual message parts) cannot change, ever.
Each message also has mutable attributes, the most important being {\em flags}.

Messages in a remote IMAP mailbox are identified by two partially independent
numbers, the {\em sequence number} and the {\em UID}.  Sequence numbers start at
one and are consecutive, that is, they dynamically change when a message in the
middle of the mailbox is removed.  This means that one particular message can
have multiple sequence numbers assigned over time and that one sequence number
can refer to totally unrelated messages, as the numbering changes.

On the other hand, the {\em UIDs} work differently.  When a message is delivered
to the mailbox, it receives a UID one greater than the previous one.  After the
UID is assigned, it will never change for that particular message, and the
mailbox remembers the biggest assigned UID so far.  This means that UIDs are
never recycled over the normal life span of a mailbox and they can be used as a
persistent identifier to one message in a particular mailbox.

Sometimes a non-standard situation might happen, e.g. when a third-party
accessed the mailbox data behind the back of the IMAP server.  Perhaps the user
tried a utility program that removes huge attachment from an otherwise valuable
message, or a spam classifier modified message headers directly, bypassing
standard mechanisms for proper flagging.  In any such case, the IMAP guarantee
of immutable nature of the messages does not hold, and a compliant IMAP server
has to communicate this problem to all clients.  The {\em UIDVALIDITY} mailbox
attribute serves this purpose.  Whenever the UIDVALIDITY changes, a compliant
client is required to throw away any cached information about a mailbox because
the UID assignment starts over again.

UID and UIDVALIDITY combined together form a 64bit integer which is absolutely
guaranteed to be an unique identifier for any message in the given mailbox -- if
the IMAP client once retrieved any immutable part from mailbox A with this 64bit
identifier, subsequent queries for the part identified by the same 64bit number
in the same mailbox can be satisfied from the client's cache.

It should be noted that this UID/UIDNEXT numbering is strictly per-mailbox and
not per-server, and that different users accessing the same server might see the
same 64bit number for totally unrelated messages.  While RFC3501 \cite{rfc-imap}
strongly recommends that UIDVALIDITY should be kept constant if at all possible,
a compliant IMAP client has to deal with UIDVALIDITY changes nonetheless.

The \label{imap-envelope}{\em IMAP envelope} is a data structure holding
interesting values determined mostly from various message headers.  It is {\em
not} the same thing as the SMTP envelope, which typically stores less
information.

The {\em Internal date} holds a time stamp about when the message was delivered
to the IMAP server.  {\em Size} attribute is the number of octets required for
storing a full copy of message in RFC2822 format.  {\em Body structure}
represents the tree-like structure of the MIME message.

The \label{imap-msg-flags}{\em IMAP flags} is an attribute holding a set of zero
or more named tokens associated with the message.  A flag can be either {\em
permanent} or {\em session-only}.  Permanent flags are stored in the mailbox
itself, while the session-only flags disappear on a subsequent reconnect to the
mailbox.  Some of the flags are defined by the RFC 3501 standard itself, others
are added by various recommendations and extension standards.  An IMAP server
might allow IMAP clients to store their own message flags.  A special means of
communicating this via the protocol exist\footnote{The {\tt PERMANENTFLAGS}
response code}.

Finally, message contents itself is distributed in form of {\em message texts}.
Clients are free to choose from retrieving MIME message parts separately or as
one blob.

\section{MIME}

Multipurpose Internet Mail Extensions \cite{rfc-2045} \cite{rfc-2046}
\cite{rfc-2047}, or just MIME, is a set of standards for extending the old SMTP
mail from 7bit ASCII texts to rich multimedia contents.  Parts of the standard
deal with encapsulation of non-ASCII characters in message headers, as well as
providing instructions about how to embed non-textual data in e-mail messages.

\subsection{Message as a Container}

An important aspect in MIME is its introduction of an internal structure to
e-mail message bodies.  Previously, a message body was just a plain text, while
from the codification of this standard, it became possible to exchange more
structured information.  Contemporary users are probably familiar with HTML
e-mails containing embedded images and a PDF attachment for seamless printing.
Users routinely send e-mails with attached photographs from their holidays, or an
archive containing business presentation.

MIME standard achieves this by defining a {\em content type} for message
bodies.  This content type, usually set by an RFC 2822 header of a
message, determines how a compliant mail client displays the message.  The
old-school non-MIME mails have a default MIME type of {\tt text/plain}, an
unformatted text in 7bit US-ASCII encoding.

Another portion of the standard defines generic-purpose containers, i.e.
abstract content types which contains other message parts within.  Using
these containers, a message can suddenly evolve from an unstructured blob of
text to a deep tree of unrelated body parts.

Most of these containers are defined in RFC 2045 \cite{rfc-2045} and 2046
\cite{rfc-2046}, one was added in RFC 2387 \cite{rfc-multipart-related}.

The {\tt multipart/alternative} encoding provides a mean of transmitting a
message that comes in several different formats, but each of them provides the same data.  The
MUA is expected to pick one of the underlying parts and display it, perhaps
allowing the user to override this selection as a nice bonus.

The {\tt multipart/mixed} is a generic catch-all MIME type for multipart
messages.  If a compliant MUA sees a multipart content type that it can not
recognize, it should behave as if it was the mixed multipart.  Typical action is
showing all embedded parts next to each other.

There are several more content types defined, but describing them in detail is
out of scope for this thesis.

\subsection{International Characters in Messages}

Original SMTP mail as defined by RFC 822 is suitable only for transferring
7bit data.  Nowadays, 8bit transport channels are far more common, both for representation of
international characters and transfer of binary data.  Therefore, the MIME
family defines {\em transfer encodings} for conversion of generic 8bit binary
data into a 7bit character stream and a portable way of expressing national
characters using 7bit ASCII characters only.

For the former, two standard encodings are available, the {\em Quoted printable}
and {\em Base64} (as well as a ``fake'' {\em binary} or {\em 8bit} encoding for
transport of binary data over less restrictive tunnels), the latter is supported
by a similar approach, as defined in RFC 2047 \cite{rfc-2047}.

\subsection{MIME Support in IMAP}

A large part of work involving the MIME support can be performed by the IMAP
server.  This is especially true for message structure parsing, where the IMAP
specification provisions ways to explore the ``tree'' of a message, as well as
methods for retrieving the resulting body parts separately (and even including
byte ranges for more convenient download).

The IMAP protocol does not, however, handle the rest of MIME support, nor much of
other header parsing.  While it is common to ask the server for just a named
subset of RFC 2822 mail headers, the returned data are not pre-parsed in any
way.  Therefore, much of the work from coalescing several header parts to
decoding international character data is left as an exercise to the client.

Fortunately, various message attributes in the IMAP protocol can serve as a
substitute for parsing the RFC 2822 headers.  Sender and recipient data is
available from the IMAP {\tt ENVELOPE}, message size as {\tt RFC822.SIZE}, IMAP
itself provides extensions for message threading \cite{rfc-threading} et cetera.

\section{IMAP Protocol Flow}

Upon a successful connection to a remote IMAP server, client might choose to
authenticate itself if the server has not already pre-authenticated it automatically
under a particular user account.  After that or in case the authentication is
not required, the connection enters the {\em authenticated} state.  In this
phase, no mailbox is selected and only a subset of commands is valid.  Clients
can, for example, ask for listings of the mailbox tree, get quick information
about number of messages in a particular mailbox or otherwise manage mailboxes
as a whole.  They have to, however, {\em select} a mailbox in order to do
anything else, like retrieving messages or marking them as read.

A mailbox can be opened as read-only or for both reading and writing, provided the
authenticated user has sufficient privileges.  Retrieving messages can be done in
both modes, but write operations (like storing a new message to the mailbox or
manipulating the message flags) require read-write access.

\subsection{Commands and responses}

IMAP is a line-oriented text protocol.  Anything that server sends to its client
is called a {\em response} while all data flowing in the opposite direction are
known as {\em commands}.  Each command generated by client begins with a unique
identifier called {\em tag} which is followed by the actual command identifier.
Most commands also accept (or require) a {\em parameter} which might be a string
literal, a number, an atom, a parenthesized list of other types etc.
Transmission of generic string literals requires an explicit confirmation from
the server unless the LITERAL+ \cite{rfc-literalplus} extension is active.
Support for this extension is crucial for good performance, as if absent, most
commands will be delayed by at least network round-trip-time and no support for
pipelining of such commands would be possible.  Clients are free to issue several
commands without waiting for their completion, subject to certain ambiguity
rules\footnote{See section 5.5 of RFC 3501}.

Each command usually leads to at least one server response, the {\em tagged
reply}.  This response contains the same tag as the client sent in the original
command, allowing the client to find out if the command failed or succeeded.
Some commands may result in sending further replies, the {\em untagged
replies}.  An important detail of the protocol is that such untagged replies are allowed to occur
at almost any time during the conversation and {\em not} only as a result of a
particular command.  A compliant client must be prepared for this situation and handle
it correctly, for there is no difference between those unilaterally generated
and ``requested'' replies.

\subsection{Mailbox Synchronization}

For efficiency reasons, each client usually keeps some data in its persistent
cache.  Good candidates for caching are message texts, a copy of body structure or
perhaps even flags from the previous session.  Because some server replies, namely
the {\tt EXPUNGE} which informs the client that a message has been removed from
the mailbox, refer to messages only by their sequence numbers and clients
usually identify messages by UIDs in order to allow message texts caching, a
compliant client has also to keep its UID-to-sequence number mapping up-to-speed
with the server.  Doing so for just a visible (in the GUI) part of the mailbox
might seem tempting, but the risks involved therein (like not knowing what to
purge from local cache as a result of received EXPUNGE without extra {\tt
FETCH}) are high and usually outweigh possible initial savings.

An IMAP server fully conforming to RFC 3501 is required to sent a few ``status''
replies as a part of selecting of a mailbox.  Interesting fields for the mailbox
synchronization are updates on UIDVALIDITY, UIDNEXT and EXISTS.  After receiving
all three, client has enough information to decide about what actions are
necessary to bring the particular mailbox back to a fully synchronized state.
Clients only have to deal with mutable message attributes, which are the
sequence number, the UID number and message flags.

If the UIDVALIDITY value changed, client is required to throw away any data it
might have in its local cache that refers to this particular mailbox.  Some UIDs
might have changed, message parts got modified or the client just happens to
talk to a buggy server implementation.  This situation is then identical to
syncing form an initial state where client has no knowledge about the target
mailbox. In the following text, we will assume that the UIDVALIDITY value is the
same as on the previous synchronization.

If the UIDNEXT value did not change either, no messages were delivered to the
mailbox since the last time.  If the EXISTS number is the same as well, there
have been no changes at all and the only item to re-synchronize are message
flags.  If EXISTS number changed, some messages were permanently removed.
Clients usually re-fetch their UID-to-sequences mapping as a whole, but
implementations which prefer to transfer the lowest amount of data and value
that above number of iterations over the network could use an algorithm similar
to binary search to find out which regions of the UID mapping were left intact
and which require re-synchronization.

If the UIDNEXT changed\footnote{It could only grow up, because UID numbers are
strictly ascending} and the $\Delta UID$ is the same as $\Delta EXISTS$, some
new messages arrived, but none were deleted and the client only has to ask for
UIDs of these new messages (and FLAGS for all messages in the mailbox).  If the
increments of UIDNEXT and EXISTS differ, a fallback to generic
synchronizing approach where client retrieves UIDs of all messages in the mailbox is
necessary.  It is also reasonable to transfer flags as a part of this step.

If there were no cached information or a UIDVALIDITY change occurred, clients
would typically fetch a complete UID mapping and message flags.

After performing these steps and having received all relevant replies, the
client is in a fully synchronized state and ready for receiving mailbox updates.

\subsection{Changes to Mailbox}

The mailbox that is selected is typically of some interest to the user.  The
original IMAP RFC mandated that all clients must be prepared for receiving
updates about number of messages contained therein, removal of old ones and
deliver of new, but due to the lack of support by clients at that time, this feature
was rarely fully utilized.  Instead, the {\em NOOP} command was widely used as a
polling check for new messages.  However, this had all the disadvantages about
polling-based approach, especially unneeded increase in network traffic, lower
possibility to enter deep sleep for battery powered devices and even delays in
mail delivery, for some clients were rather hesitant to poll for changes too
often.  A solution to this problem is presented in RFC 2177 \cite{rfc-idle}
where client enters a special mode (which is syntactically similar to
synchronizing literals that are used when support for LITERAL+
\cite{rfc-literalplus} is absent) to tell the server that it can really send
real-time notifications to the client.

Changes to number of messages in the mailbox are communicated mainly by two
kinds of untagged replies, {\em EXISTS} and {\em EXPUNGE}.  EXPUNGE reply
contains a sequence number of message that has been permanently removed from the
mailbox and immediately causes all messages with sequence number higher than
the message being deleted to decrement their assigned sequence numbers by one.
This means that clients have to work with sequence numbers, unless they are
prepared to update their UID map on each expunge, which would be a rather
expensive approach.  EXISTS reply is used to inform the client that new message
was delivered, or it could follow an EXPUNGE to provide a redundant information
to the client about number of messages in the mail store.  Under no
circumstances it could be used to decrease the number of messages in current
mailbox, as this is already handled by EXPUNGE.

\subsection{Fetching and Manipulating Messages}

Once selected, messages in the mailbox can be queried and downloaded to the
client.  The {\em FETCH} command serves for retrieving all data from body
structures and IMAP envelopes to actual message parts.  Message flags are
updated by the {\em STORE} command and its variants.  Existing messages can be
copied to other mailboxes with the {\em COPY} command, new messages can be
stored into a mail folder by command {\em APPEND}.

\subsection{Queries Against Other Mailboxes}

Because most graphical MUAs want to show the user more than one mailbox at a
time, at least in a collapsed view with counts of new or unread messages, and
frequent switching among several mailboxes is not a cheap operation, IMAP also
provides a command for determining state of other than currently selected mailbox,
the {\em STATUS} command.  Using this command for querying state of a currently
selected mailbox is explicitly forbidden by the RFC, for a conforming client
already has all required information to compute all possible message counts.
Another reason for disallowing STATUS on current mailbox is that in some server
implementations, this operation might cause a process to open the mailbox once
more which might in turn provide inconsistent results and confuse the client
even more.

\subsection{Searching, Sorting and Threading}

Even the base specification of the IMAP protocol provides facilities for
searching the messages stored in current mailbox.  This was extended over time
by ways to request specific message ordering or threading \cite{rfc-threading}.
Most e-mail clients therefore do not need to have all message bodies and headers
fetched and available in local cache in order to allow their users to search
through messages.  However, this searching and sorting is not available in
offline clients which have to implement their own code for such purposes, or
disallow these operations when network access is disabled.

Following revisions also specified certain extensions related to
internationalization of e-mail traffic.  For example, searching for non-ASCII
characters in messages would be very hard using the bare-bone SEARCH command
only, as each message in mailbox can be store in one of the myriad encodings
available in MIME standard.  Subsequent revisions to IMAP therefore defined
support for more intelligent searches which support unified searching through
message that arrived in different character sets.

\subsection{Manipulating Mailboxes}

Due to some historic reasons (especially the old versions of SMTP protocol being
7bit), IMAP does not use any standard encoding for mailbox names.  Instead, a
modified version of UTF-7 called simply {\em modified UTF-7} is mandated.  The
standard does not specify what to do if the other side uses 8bit data in mailbox
name, but most implementations nowadays fall back to UTF-8, which is arguably a
good thing to do.

Mailboxes can be created, deleted and renamed by IMAP clients.  If the server
supports it, they form a tree hierarchy.  Some servers distinguish between
mailboxes for storing messages and those that contain only other mailboxes,
while other implementations allow free mixing of both.  Unfortunately there is
no reliable way to tell to which class the target server belongs, but provisions
are given about how to create a folder for message and how for other mailboxes.
If mailboxes form a tree, the server implementation specifies which character is
used as a hierarchy delimiter.  Most common choices in today's mainstream IMAP
servers are dot ({\tt .}) which is common among servers exporting Maildir
folders and slash ({\tt /}) as used by servers exporting traditional mbox UNIX
mail store.  This inconsistence might lead to user confusion, but cannot be
avoided without breaking backward compatibility.

\subsection{IMAP Extensions}

FIXME.

\section{Other Methods of Mail Store Access}

IMAP is not the only protocol for accessing mails stored on a remote server.  In
this section, we will introduce some of the alternatives and discuss their
advantages and disadvantages when compared to IMAP.

\subsection{POP3}

The main difference between IMAP and POP3\footnote{Post Office Protocol version
3} protocol is their philosophy --- IMAP was designed for scenarios where user
leaves all her e-mail at a central place which has authoritative answer (FIXME answer?) about
what messages are supposed to ``be there'' while POP3 expects their users to
download all e-mail to a computer where the POP3 client runs and then {\em
permanently remove the original messages from the server}.  Upon each reconnect,
POP3 client downloads and deletes all new messages again.  This fundamental
difference leads to troubles when user wants to use POP3 in situations it was not
designed for -- a classic example is accessing one mailbox from two places over
time, like a work computer and a laptop at home.  While there are certain ways
to work around these POP3 design limitations which allow leaving a copy of each
message on the POP3 server, they all remain just hot fixes
instead of providing a proper and reliable solution.  An excellent comparison of
these two paradigms is available from Grey's paper \cite{imap-vs-pop}.

\subsection{MAPI}

A proprietary protocol from Microsoft which documentation got recently opened
for public access after a pressure from both government an non-government
organisations.  It has some nice features comparable to IMAP, but offers no real
advantage besides tight integration with other Exchange services (which might or
actually might not be considered an advantage).  Implementations from other
parties than the original vendor are scarce.

\subsection{Webmail}

Several years ago, a boom of online services from well-known companies like
Google or regional Internet portals grew in popularity.  Despite inherent
limitations of web-based user interfaces, most ``casual'' users of e-mail,
especially the second ``home mailboxes'' of regular employees, are today hosted
and accessed via a webmail interface.  Although hardly any of these services
offers any guarantee about availability or performing backups, regular users
typically are not concerned about these issues.  Reliability of these services
is, on the other hand, typically greater than that of a typical home PC with
aging hardware, and friendliness of the UI improved with advent of recent
web technologies.

\subsection{IMAP Criticism}

During more than 20 years of the IMAP history, the standard has certainly
accumulated some aging dust.  A typical example are the synchronizing literals,
originally the only way to send more complex data to the server, which require
explicit server's acknowledgment before continuation.  Given a typical GPRS
connection which is still common in many countries, a latency of 500ms can delay
all interaction with the IMAP server considerably.  Many similar aspects were
fixed by subsequent modifications to the standard, others are rather unfixable
design limitations.

A fine example of such limitation, at least when perceived by contemporary
optics, is the concept of mailboxes.  Online messaging services like GMail
introduced a new way of thinking about messages -- previously a message would
belong to exactly one mailbox and if user wanted to have it available from two
places (perhaps in ``work related'' and ``mailing list'' folders), she would
have to store two copies.  IMAP has no concept of a ``link'' between two
messages.  When one is flagged, the other copy is left intact.  Another
prehistoric feature of IMAP is its modified UTF-7 encoding which is not used
anywhere else besides the IMAP protocol.  Another point -- IMAP usually does a
good job in shielding the client implementation from the need to parse all MIME
data, but certain pieces remain, like the RFC 2822 headers with international
characters.

Some authors also criticize the protocol on basis of its complexity and being
loosely-defined in certain aspects.  An example is the IDLE \cite{rfc-idle}
extension, which adds an explicit mode indicating to the server that the client
is really ready to process any updates to mailbox state, while the original RFC
mandated the client to be in such condition all the time.  Most common MUAs also
have troubles with getting all side-corners of IMAP right.

That being said, the IMAP is the only widely implemented standard for accessing
an online or disconnected mail store.  It has certain deficiencies, but all in
all, it is a perfectly usable protocol.

\chapter{\trojita Design}

\section{Overview}

\section{Model-View Architecture}

\section{Parser}

\section{Low-level Parser}

\section{Cache}

\chapter{Other IMAP Implementations}

\chapter{User's Guide}

\section{About Trojita}

\trojita is a standards-compliant IMAP Mail User Agent. Designed with portability
and efficiency in mind, it allows the user to operate over congested network
lines with high latency, while still providing sufficient comfort.

\section{Installation}

\subsection{Prerequisites}

In order to build \trojita from source, some libraries and helper programs are
required to be present on the system.

\begin{enumerate}
    \item{CMake 2.6}
    \item{Qt 4.4}
    \item{Working implementation of the TR1 C++ standard}
\end{enumerate}

Recommended version of the Qt framework \cite{qt} is at least 4.5. \trojita,
however, includes compatibility code for older versions of Qt. Minimal supported
version is Qt/4.4. Any older release cannot be supported, as the WebKit HTML
engine which we use for e-mail rendering rendering was added during the 4.4
development. If users choose to build against 4.4, some \trojita features might
be missing.

Unlike most Qt projects, \trojita uses CMake \cite{cmake} instead of {\tt qmake},
especially due to author's familiarity with the former system.

\trojita makes use of some standard C++ features that were adding during the TR1
\cite{std-tr1} standardisation, namely the {\tt std::tr1::shared\_ptr} smart
pointer. A working implementation of said class is required.  \trojita build
system will make use of compiler's built-in TR1 support, if available, and
resort to using Boost if the compiler does not support this extension.

Some 3rd party code is also bundled in the source code distribution.  A list of
such products including the reason for not just using an external library
follows:

\begin{enumerate}
    \item{QwwSmtpClient \cite{qwwsmtpclient} is a Qt SMTP client library written
        by Witold Wysota.  Released to public only in mid May 2009, obtaining a
        package of it might prove challenging.  The copy in \trojita's codebase
        also contains some fixes which did not made it to the publicly released
        version yet.}
    \item{ModelTest \cite{modeltest}, a class for verifying certain assumptions
        of the Interview Model-View framework of Qt.  Qt Software (former
        Trolltech) does not ship it with Qt sources, so \trojita includes a copy
        for developers' convenience.  This feature is not turned on by default
        in release builds.}
    \item{QtIconLoader \cite{qticonloader}, another helper library from Qt
        Software.  Purpose of it is providing a unified icon set, as per
        relevant XDG standards.}
    \item{KCodecs \cite{kcodecs}, an implementation of UTF-7, Base64,
        Quoted-Printable and related RFC 2047 standards.  These functions are
        reasonably self-contained, so they are bundled with \trojita instead of
        depending on quite heavyweight {\tt kdelibs}.}
\end{enumerate}

\subsection{Obtaining source code}

Source code is available from the CD attached to this thesis or online at FIXME.

\subsection{Installing}

Unpack the source tarball, prepare the build environment and launch the compile
process by entering these commands:

\begin{enumerate}
    \item{{\tt tar -xf trojita-1.0.tar.bz2}}
    \item{{\tt cd trojita-1.0}}
    \item{{\tt mkdir \_build}}
    \item{{\tt cd \_build}}
    \item{{\tt cmake -DCMAKE\_BUILD\_TYPE=RelWithDebInfo ..}} (the {\tt ..} are
        not a typo)
    \item{{\tt make}}
\end{enumerate}

After a while, a new binary, {\tt trojita}, will appear in the build directory.
Launch it to start \trojita.

\section{Using \trojita}
\subsection{Configuration}

All configuration is, thanks to the Qt's QSettings, stored in target platform's
common storage location. The only supported mean of configuring \trojita is
through its built-in configuration dialog that will automatically appear on the
first run.  It can also be launched any time later from the application menu.

% FIXME: screenshot

On the first page, user is asked to provide identity data. These are used for
composing e-mail messages.

Second page is the most important one in configuration of \trojita. It is the
place where the user tells the application which IMAP server is it supposed to
talk to, which credentials should be used for authentication et cetera.

One unusual thing which is not common among other MUAs is the choice of how to
connect to the IMAP server.  Most users are probably familiar with connections
over TCP, optionally secured by SSL/TLS, but the last option, the {\em Local
Process}, is unique to \trojita. If selected, it replaces the usual TCP
connection to the IMAP server with launching a local process, which is in turn
used instead of the remote IMAP server. The most interesting use case is to set
up SSH connection to the remote server (along with proper SSH keys) and use that
as a secure, reliable and convenient way for accessing e-mail, without having to
enter passwords all the time. For example, to launch an instance of Dovecot IMAP
server on host {\tt mail.example.org}, user should enter {\tt ssh
mail.example.org dovecot --exec-mail imap}.

The last page serves for configuring the message composition and sending.
\trojita supports mail submission via SMTP\footnote{Simple Mail Transfer
Protocol} and local {\tt sendmail}-compatible MSA\footnote{Mail Submission
Agent\cite{rfc-msa}, a local gateway whose purpose is forwarding e-mail to the
network of MTAs} application.

\subsection{Basic usage}

After the configuration is complete, the user is presented with a familiar user
interface. The left pane presents a view of remote mailboxes, or folders, on the
IMAP server. The top widget contains list of messages in currently opened
folder, and finally the bottom area is reserved for displaying of the actual
e-mail messages.

Depending on the network parameters and server load, it might take a while
before the list of mailboxes appears. The delay is only significant when
launching \trojita for the first time, for the retrieved list of mailboxes is
cached on a persistent disk cache for later use.

After the list of mailboxes is loaded, user can click on them to open them in
the {\em Message List} widget on the bottom right corner of the application
window.  After having selected a mailbox, the list of messages appears shortly.
Again, the loading is faster next time \trojita is launched.

Depending on user settings, \trojita typically will not transfer all messages
metadata upon opening a mailbox.  This feature is implemented to conserve
both server's and client's resources and to be nice to the network traffic.  In
a typical mode, that is, when not being told to be extra-mean with resources,
\trojita fetches only metadata of those messages that are displayed on the
visible screen region, and some more to make scrolling smoother.  More on this
topic is said in another chapter of this thesis.

Message contents is transfered on-demand only, that is, when user selects a
message for viewing.  As before, all persistent contents once downloaded from
the server is cached in a persistent on-disk cache for future use.  This applies
to both message metadata and the real contents of e-mails.

\subsection{Managing Messages and Mailboxes}

Messages displayed to user are automatically marked as read after some delay.
The idea behind this feature is to allow users to peek on a particular message,
shortly judge its contents, classify it as a ``process later'' and then move on
to another one.  This is a common approach shared with other MUAs as well.

Manipulating with {\em message flags} (see chapter \ref{imap-msg-flags}) is possible
via the context menu (appearing on the message list) as well by standard
keyboard shortcuts.  The {\tt M} key toggles the {\em Read} status and {\tt
Delete} key schedules messages for removal.

In consistency with usual IMAP approach (and unlike the common modus operandi on
mainstream desktop MUAs), \trojita has no concept of {\em Trash}
\footnote{A {\em Trash} is a mailbox that presents a view of all deleted
messages.  It might or might not be a real mailbox on the server, or just a
virtual folder implemented purely in the e-mail client.} for
messages.  Instead, each message can be marked as deleted.  If flagged as such,
\trojita displays a graphical indicator next to the message subject.  A deleted
message will not be physically removed from the mailbox until the user explicitly
requests such an action.  This operation, called {\em expunge} in the IMAP
dictionary, is available from the application menu.  Expunging happens on one
mailbox at a time, there is no way to tell the IMAP server to permanently remove
all deleted messages from all mailboxes.

To create new mailbox, user is expected to use the context menu on the left {\em
mailbox view} pane.  Due to technical limitations, some IMAP servers distinguish
among folders that can contain other mailboxes and folders for message store
only.  Additionally, there is no reliable way for the IMAP client to tell
whether the remote server is affected by this limitation or not.  Therefore,
\trojita uses the same approach as other common MUAs and ask user what kind of
mailbox she wants to create.

% FIXME screenshot

Mailbox removal is triggered from the context menu as well.  Due to the possibly
serious consequences of this action (there is no way to undelete a removed
mailbox), user is presented with a confirmation dialogue asking her whether she
is serious in her intent to remove the mailbox.

Messages can be copied or moved between mailboxes by simple drag-and-drop
operations.  To place a copy of one or more messages into another mailbox,
simply select the messages to copy in the opened mailbox and then drag them with
mouse to the target mailbox.  If user wants to move them instead of performing a
copy, she is supposed to hold the {\tt Shift} key.  This behavior is determined
by the Qt framework, so it might vary among the supported platforms.

\subsection{Message Composition}

Message Composition widget is triggered by the corresponding item in the menus.
The displayed dialog should be familiar to most users.

\subsection{Checking Mail}

\trojita automatically checks the current folder for new messages.  If the
remote server supports the IDLE extension \cite{rfc-idle}, the method used is
extremely friendly to both network and server resources.  In cases the server
does not support this feature, \trojita resorts to polling with all it
disadvantages.

\section{Advanced usage}

\subsection{Partial Message Fetching}

Thanks to unique features of the IMAP protocol, it is not necessary to fetch the
whole message in order to display just a fragment of it.  \trojita handles this
automatically, fetching message parts on demand.  It also will not transfer message
metadata (like its RFC2822 envelope) unless really needed.

\subsection{Offline Mode}

When put into {\em offline mode}, \trojita will prevent itself from accessing
the network at all.  Users will be still able to view already transfered data, but
no checks for new messages will be performed and those that are not stored in
local offline cache will not have actual message bodies displayed until the user
selects either {\em online} or {\em expensive network} mode from the menu.

\subsection{Expensive Network}

Some users connect to their IMAP servers over a slow or otherwise unsuitable
network.  Examples include dial-up connections or various mobile solutions like
GPRS\footnote{General packet radio service (GPRS) is a mobile data service
typically offering low bandwidth, high latency Internet connection with billing
based on the amount of transfered data} where the overall convenience of having
all messages immediately available is usually held back by the need to conserve
network resources.  In such environments, it might be desirable to fetch only
the absolute minimum data required to provide properly synchronized view of the
remote mailbox.

In this situation, user is expected to choose the {\em Expensive Network} mode
from the application menu.  \trojita will not preload data to improve
interactivity, nor will it show all message parts by default.  Upon viewing a
message, only parts smaller than a certain size are transfered automatically,
others will be loaded after clicking a button.  This on-demand loading is
implemented separately for each part, allowing great flexibility for user to
choose what is interesting and what can be ignored.

\subsection{Online Mode}

{\em Online mode}, the common mode of operation, is what users are familiar from
other IMAP clients.  \trojita employs intelligent preloading of interesting
contents to minimize waiting times, yet does not blindly download everything,
including parts that will not be ever shown.

\subsection{Local Cache}

All data once transfered form the network\footnote{Subject to some restrictions
like maximal message part size and cache parameters} are kept in a persistent
offline on-disk cache.  If a message data is transfered twice, it means that at
least one of the following conditions was met:

\begin{enumerate}
    \item{The server could not guarantee the immutability of messages.  This might
        happen when some process touches on-disk data internal to the IMAP
        server in a non-standard way.}
    \item{User has changed her settings.  Whenever any data used to access the
        remote server changes, the only safe thing to do is to throw away
        anything in the local cache and start over again.}
\end{enumerate}

\subsection{Full Tree View}

As an optional part of the GUI, the {\em full tree view}, activated by the
corresponding item in the menu, provides a full view of the server contents,
including mailbox hierarchy, messages in mailboxes and mailbox parts.  This
widget is a nice demonstration of a complex nature of MIME messages -- while
most users are probably familiar with the concept of attachments, the truth is
that MIME defines a fully recursive structure that one message can consist of.
For example, it is perfectly possible to have a message composed of six parts,
each of them being a forwarded message, and each of these forwarded messages
being rich-structured as well.  \trojita does its best at formatting the message
in a sane and familiar way, but unfortunately some senders do not really follow
the same philosophy and produce messages whose structure leaves much to be
desired.  More on this matter can be found in another place in this thesis
\ref{FIXME}.

\chapter{Conclusion}

\trojita, the IMAP e-mail client implemented for this thesis,\ldots

\begin{thebibliography}{99}
\addcontentsline{toc}{chapter}{Bibliography}
    \bibitem{rfc-imap}Crispin, M.: Internet Message Access Protocol - Version
        4rev1, {\em RFC3501}, March 2003
    \bibitem{rfc-literalplus}Myers, J.: IMAP4 non-synchronizing literals, {\em RFC
        2088}, January 1997
    \bibitem{rfc-idle}Leiba, B.: IMAP4 IDLE command, {\em RFC2177}, June 1997
    \bibitem{crispin-ten-commandments}Crispin, M.: Ten Commandments of How to
        Write an IMAP client, September 2006, Washington,
        {\tt http://www.washington.edu/imap/documentation/commndmt.txt.html}
    \bibitem{client-best-practices}Sirainen, T., Cridland, D. et al.: Best
        Practices for Implementing an IMAP Client, {\tt
        http://www.imapwiki.org/ClientImplementation}
    \bibitem{rfc-2822}Resnick, P. (editor) et al.: Internet Message Format, {\em
        RFC2822}, April 2001
    \bibitem{rfc-msa}Gellens, R., Klensin, J.: Message Submission, {\em
        RFC2476}, December 1998
    \bibitem{rfc-2045}Freed, N., Borenstein, N.: Multipurpose Internet
        Mail Extensions (MIME) Part One: Format of Internet Message Bodies, {\em
        RFC2045}, November 1996
    \bibitem{rfc-2046}Freed, N., Borenstein, N..: Multipurpose Internet
        Mail Extensions (MIME) Part Two: Media Types, {\em RFC2046}, November 1996
    \bibitem{rfc-2047}Moore, K.: MIME (Multipurpose Internet Mail Extensions)
        Part Three: Message Header Extensions for Non-ASCII Text, {\em RFC2047},
        November 1996
    \bibitem{rfc-multipart-related}Levinson, E.: The MIME Multipart/Related
        Content-type, {\em RFC 2387}, August 1998
    \bibitem{rfc-content-disposition}Troost, R., Dorner, S., Moore, K.:
        Communicating Presentation Information in Internet Messages: The
        Content-Disposition Header Field, {\em RFC2183}, August 1997
    \bibitem{rfc-threading}Crispin, M., Murchison, K.: Internet Message Access
        Protocol - SORT and THREAD Extensions, {\em RFC5256}, June 2008
    \bibitem{imap-vs-pop}Gray, T.: Message Access Paradigms and Protocols,
        September 1995, online at {\tt
        ftp://ftp.cac.washington.edu/mail/imap.vs.pop}
    \bibitem{qt} Qt Reference Documentation, {\tt http://doc.trolltech.com/4.5/}
    \bibitem{cmake} CMake Documentation, {\tt
        http://www.cmake.org/cmake/help/documentation.html}
    \bibitem{std-tr1} ISO/IEC TR 19768, C++ Library Extensions
    \bibitem{qwwsmtpclient} Wysota, W.: QwwSmtpClient released, May 2009, \\
        {\tt
        http://blog.wysota.eu.org/index.php/2009/05/13/qwwsmtpclient-released/}
    \bibitem{modeltest} Qt Software, ModelTest, February 2008, {\tt
        http://labs.trolltech.com/page/Projects/Itemview/Modeltest}
    \bibitem{qticonloader} Qt Software, February 2009, \\ {\tt
        http://labs.trolltech.com/blogs/2009/02/13/freedesktop-icons-in-qt/}
    \bibitem{kcodecs} The KDE Project: kdelibs, {\tt
        http://api.kde.org/4.x-api/kdelibs-apidocs/}
\end{thebibliography}

\end{document}
